\section{The excess measured by the determinant}\label{sec:excess}
Fix a permissible pair $(q_0,d)$ as in Hypothesis~\ref{hyp:scale}. Thus $q_0$ is a positive multiple of $18480$, and $d$ is a prime congruent to $q_0/2-1\pmod {q_0}$. In particular $d\equiv7\pmod 8$ and $(q_0,d)=1$.
For a modulus $q\ge1$ coprime to $d$, write $\chi_d(n)=(n/d)$ and define
\begin{align}
 \lambda_{d,q}(n)&=\mathbf 1_{n\equiv1\bmod q}\Lambda(n)(1-\chi_d(n)),\notag\\
 R_{d,q}(y)&=\sum_{1\le n\le y}\lambda_{d,q}(n)(1/n-1/y)
             -\frac{\log y}{\varphi(q)}-c_{d,q},\qquad y\ge1.\label{eq:R}
\end{align}
Take $c_{d,q}=G_{d,q}(1)$, where $G_{d,q}$ is the logarithmic derivative average in \eqref{eq:p1G}. This value is real and is the unique choice for which $R_{d,q}(y)\to0$ (Lemma~\ref{lem:pnt}). For $1\le z\le a$, put $v_z(n)=(a/n-a/z)_+$. Then
\begin{equation}
 \sum_{n\le a}\lambda_{d,q}(n)v_z(n)
 =a\left[R_{d,q}(z)+\frac{\log z}{\varphi(q)}+c_{d,q}\right].
 \label{eq:scaled}
\end{equation}
Here $w_+=\max(w,0)$, and sums in $n$ start at $1$.

\subsection{Fixing the centering constant}
Let $Q\ge1$ and let $d$ be an odd prime coprime to $Q$. Write $L_\eta(s)=L(s,\eta)$, with characters extended by zero on nonunits and the Euler product taken at the indicated level. By classical continuation and nonvanishing for Dirichlet $L$-functions, the only singularity in $\Re s\ge1$ is the pole at $1$ for the principal character, and there are no zeros there. To remove the pole, set
\[
 D_\eta(s)=
 \begin{cases}
 (s-1)L(s,\eta),&\eta\text{ principal},\\
 L(s,\eta),&\eta\text{ nonprincipal}.
 \end{cases}
\]
Each $D_\eta$ is entire and $D_\eta(1)\ne0$. Each twist $\eta\chi_d$ is nonprincipal at level $Qd$. By the Chinese remainder theorem it takes the value $-1$ on an integer congruent to $1$ modulo $Q$ and to a quadratic nonresidue modulo $d$.


With the weights in \eqref{eq:R}, put
\begin{equation}\label{eq:p1G}
 G_{d,Q}(s)=\frac1{\varphi(Q)}\left[
  \sum_{\eta\bmod Q}\left(\frac{D'_{\eta\chi_d}}{D_{\eta\chi_d}}
                             -\frac{D'_\eta}{D_\eta}\right)(s)-1\right].
\end{equation}
For $\Re s>1$, orthogonality and the absolutely convergent logarithmic derivatives give
\begin{align}
 \sum_{n\ge1}\lambda_{d,Q}(n)n^{-s}
 &=\frac1{\varphi(Q)}\sum_{\eta\bmod Q}
       \left(\frac{L'_{\eta\chi_d}}{L_{\eta\chi_d}}-
             \frac{L'_\eta}{L_\eta}\right)(s)\notag\\
 &=G_{d,Q}(s)+\frac{s}{\varphi(Q)(s-1)}.\label{eq:p1Dir}
\end{align}
For the principal character, $D'_1/D_1=L'_1/L_1+1/(s-1)$. The constant $-1$ in \eqref{eq:p1G} supplies the remaining term in $s/(s-1)$. All logarithmic derivatives in \eqref{eq:p1G} are regular at $1$, and \eqref{eq:p1Dir} is real on $(1,\infty)$. Passing to the limit gives $c_{d,Q}=\Re G_{d,Q}(1)=G_{d,Q}(1)$.


\subsection{Counting the places}
In an admissible datum the CM field is $E=\QQ(\sqrt{-d},\zeta_{q_1})$, where $q_0\mid q_1$ and $q_0,q_1$ have the same prime divisors. The infinite places of $E$ are identified with $\Phi=(\mathbb Z/q_1)^\times$ by choosing the representatives extending a fixed embedding of $\QQ(\sqrt{-d})$. Denote by $\Phi_a$ the full fibre of $1\bmod q_0$ and by $\Phi_b$ the prescribed sparse subset in the fibre of $-1\bmod q_0$. Write
\[
 A=|\Phi_a|,\quad B=|\Phi_b|,\quad f=|\Phi|=A\varphi(q_0),\quad
 L=\log b,\quad A\le BL.
\]
With $x\le x_1\le a=xb$, the weight $\kappa(\tau)$ equals $b=a/x$ on $\Phi_b$ and $a/x_1$ elsewhere. For an integer $t$ let $u(t)$ be its class if it is a unit modulo $q_1$, and the identity of $\Phi$ otherwise. For $0\le y\le a$ put
\begin{align}
 S(y,\tau)&=\sum_{n\le a}\lambda_{d,1}(n)
       \big(y/n-\kappa(u(-n)\tau)\big)_+,\notag\\
 S_0&=a\left(\log a-\frac1f\sum_{\tau\in\Phi}\log\kappa(\tau)+c_{d,1}\right),
 &\Tot_D&=\sum_{\tau\in\Phi_a}\big(S(a,\tau)-S_0\big).\label{eq:total}
\end{align}

\begin{proposition}[Two-cutoff identity]\label{prop:count}
For every admissible datum,
\begin{equation}
 \Tot_D=a\left\{B\big[R_{d,q_0}(x)-R_{d,q_0}(x_1)\big]+A R_{d,1}(x_1)\right\}.
 \label{eq:count}
\end{equation}
\end{proposition}
\begin{proof}
Fix $n$. If $n\equiv1\bmod q_0$, it is a unit modulo $q_1$ and multiplication by $u(-n)$ maps $\Phi_a$ bijectively to the fibre of $-1$. Exactly $B$ images belong to $\Phi_b$. If $n$ is a unit in any other class, the image fibre misses $\Phi_b$. If $n$ is a nonunit, $u(-n)=1$ and $\Phi_a$ is disjoint from $\Phi_b$. Hence the sum over places of $S(a,\tau)$ is
\[
 \sum_{n\le a}\left[A\lambda_{d,1}(n)v_{x_1}(n)
    +B\lambda_{d,q_0}(n)\big(v_x(n)-v_{x_1}(n)\big)\right].
\]
Apply \eqref{eq:scaled} at the two cutoffs. The two copies of $c_{d,q_0}$ cancel. The remaining main terms sum to
\[
 a\left[A(\log x_1+c_{d,1})
       +\frac{B}{\varphi(q_0)}\log(x/x_1)\right].
\]
This equals $AS_0$ since $\sum_{\tau}\log\kappa(\tau)=f\log(a/x_1)-B\log(x/x_1)$ and $f=A\varphi(q_0)$. Subtracting $AS_0$ gives \eqref{eq:count}.
\end{proof}

\begin{corollary}[One datum]\label{cor:one}
For each admissible datum $D$ satisfying the pointwise estimate \eqref{eq:det}, one has
\begin{equation}
 R_{d,q_0}(x)\le\frac{\varphi(q_0)L^6}{m}
               +|R_{d,q_0}(x_1)|+L|R_{d,1}(x_1)|.\label{eq:one}
\end{equation}
\end{corollary}
\begin{proof}
Solve \eqref{eq:count} for $R_{d,q_0}(x)$ to obtain
$\Tot_D/(aB)+R_{d,q_0}(x_1)-(A/B)R_{d,1}(x_1)$. Since $aB>0$, \eqref{eq:det} bounds the first term from above by $(A/B)\varphi(q_0)L^5/m$. Use $A/B\le L$ and take absolute values of the two secondary-cutoff errors.
\end{proof}

The saving $1/m$ at the primary cutoff comes from \eqref{eq:det}. Section~\ref{sec:comparison} derives \eqref{eq:det} under the arithmetic hypotheses by comparing determinants on the same family of lifts.

\section{Admissible data at every large scale}\label{sec:data-supply}

Fix $q_0>0$ divisible by $18480$ and a prime
$d\equiv q_0/2-1\pmod {q_0}$. Thus $d\equiv7\pmod8$ and
$(q_0,d)=1$. Put $K=\mathbb Q(\sqrt{-d})$.
We construct the ranks, place weights, and levels of an admissible datum for each sufficiently large prescribed length $x$.

\begin{theorem}\label{thm:p1-data}
Fix $0<\eta<1/100$. Every sufficiently large integer $x$ occurs as the
length of an admissible datum for $(q_0,d)$. With $L=\log b$, one may take
\begin{equation}\label{eq:p1-data-scale}
 x_1=\lfloor\exp(x^{\eta/6})\rfloor,\qquad
 x^\eta<L\le x^\eta+1,\qquad \sqrt{x}\le m<2\sqrt{x}.
\end{equation}
The place counts satisfy $f=A\varphi(q_0)$ and $0<A/B\le L$.
In particular $\log x_1\ge x^{\eta/12}$ eventually. The threshold for $x$
may depend on $q_0,d,\eta$.
\end{theorem}

\paragraph{Ranks, moduli and weights.}
Write $Y=\log x$, $W=x^{\eta/6}$. Choose $m=2^u$ and $r=3^v$ to be the
least powers at least $\sqrt{x}$ and $x^5$, respectively. Thus
$r<3x^5$. Set
\[
 s=22u,\quad o=2^s=m^{22},\quad
 \ell_0=3^{\lfloor\log_3s\rfloor+4},\quad
 A=3\ell_0 o,\quad q_1=q_0 A,\quad q_2=(q_0d)^{s+6}.
\]
Here $27s<\ell_0\le81s$ and $s<22+11Y/\log 2\le32Y$ eventually, so $o,q_1,q_2^2$ are
polynomially bounded in $x$ with $(q_0,d)$ fixed. Take
$E=K(\zeta_{q_1})$. Its infinite places, represented by embeddings
extending one fixed embedding of $K$, are indexed by
$\Phi=(\mathbb Z/q_1)^\times$. Indeed, $K$ has conductor $d$ and the
cyclotomic conductors $d,q_1$ are coprime. Reduction modulo $q_0$ is
surjective and all its fibres have size $A$, since $q_0$ and $q_1$
have the same prime divisors. The places of class $1\bmod q_0$ form
$\Phi_a$, of size $A$, and $f=|\Phi|=A\varphi(q_0)$.
See \src{New/Reduction/Data/FieldsAndPlaces}{the tower and its infinite places}.

For $x_1=\lfloor e^W\rfloor$, let $b$ be the least multiple of $q_1x_1$
strictly greater than $e^{W^6}$, and set $a=xb$. Since
$\log(q_1x_1)\le W+O_{q_0,d}(\log x)$, the rounding costs at most
$e^{W^6}$ for large $x$. Consequently
\[
 W-1\le\log x_1\le W,\qquad W^6< L\le W^6+1.
\]
These bounds give
$10^6Y\le\log x_1$, $10^6\log x_1\le L$, $(\log x_1)^6\le L$,
and $100\log L\le Y$. The last inequality uses $100\eta<1$.
Also $\log L$ eventually exceeds $10^{12}$ times the fixed parameter
constant in the determinant data.
Give the places in $\Phi_b$ weight $b$ and the remaining places weight $a/x_1$. Both weights are positive integers, and $a/x_1\le b$ for large $x$.

\paragraph{A large sparse fibre subset.}
Put $k=r+1-m$, $t=5r$, and $D_*=kt$. The determinant estimate requires a set $\Phi_b$ of at least $A/L$ places in the fibre of $-1\bmod q_0$ with the following property. For all primes $p\le D_*$ and integers
$0<i<j<\ell_0$, there do not exist three members $\sigma,\tau,\rho$
with $\tau/\sigma\equiv\pm p^i$ and $\rho/\sigma\equiv\pm p^j$
modulo the largest divisor of $q_1$ prime to $p$. The signs are independent
and repetitions among the three places are allowed. Write $B=|\Phi_b|$.

Here $D_*\ll x^{10}$, $o\ge x^{11}$, so $2\ell_0(D_*+1)<o$ eventually.
Choose an integer $n\ge1$, fixed in $x$ and sufficiently large that
$o,q_2^2,m\le W^n$ eventually. Since $W$ eventually dominates every
fixed power of $Y$, we also have
\[
 W\le o,\qquad(2\ell_0)^n\le W,\qquad
 3\ell_0 n2^n\le W^3 .
\]
Let $H$ be the largest power of $2$ satisfying $H^n\le o/W$.
Then $H\le W$, $o<W(2H)^n$, and $(2\ell_0H)^n\le o$.

Encode the digit vectors in $\{0,\ldots,H-1\}^n$ in base $2\ell_0H$ (cf.~\cite{behrend}).
At least $H^n/(nH^2)$ of them have the same squared Euclidean norm.
Let $S\subset[0,o)$ be the corresponding integers. Consider a congruence with positive integer coefficients $c,c'$ of the form
\[
 (c+c')e_2\equiv ce_1+c'e_3\pmod o,\qquad
 e_1,e_2,e_3\in S,\quad c,c'>0,\quad c+c'\le2\ell_0 ,
\]
where each digit on either side is at most $(c+c')(H-1)<2\ell_0H$. Thus no carries occur. Both sides are less than $(2\ell_0H)^n\le o$. The congruence is therefore an equality in each digit. The middle vector is a convex combination, with both coefficients positive, of the other two vectors on the same sphere. Strict convexity of the squared norm forces $e_1=e_2=e_3$.

Set $g=1+3\ell_0q_0$ and take $\Phi_b=\{-g^e\bmod q_1:e\in S\}$.
If $v_2(q_0)=v$, then $v\ge4$ and $g$ has order $2^s$ modulo
$2^{s+v}$ because successive squaring increases $v_2(g^{2^i}-1)$ from $v$
to $v+i$. Thus the places are distinct and in the prescribed fibre. Moreover
\[
 A=3\ell_0o
 <\frac{H^n}{nH^2}\,(3\ell_0 n2^n)WH^2
 \le B W^6\le BL .
\]
For odd $p$ the two forbidden ratios, squared to remove their signs
and raised to powers $j,i$, imply
$2j e_2\equiv2(j-i)e_1+2i e_3\pmod o$.
Hence all exponents agree. Since $2^{s+v}$ divides the part of $q_1$ prime to $p$, the first squared ratio would give $p^{2i}\equiv1\bmod 2^{s+1}$.
The identity
$v_2(p^{2i}-1)=v_2(p-1)+v_2(p+1)+v_2(i)$ for odd $p$ implies
$2^s\le (p+1)2i$, contrary to the choice of $o$. For $p=2$, every place
has the same residue modulo $3\ell_0$, so a forbidden first ratio gives
$4^i\equiv1\bmod 3\ell_0$. But $4=1+3$ has order $\ell_0$ there,
whereas $0<i<\ell_0$. This proves sparseness. These arguments specialize
the digit-sphere and order calculations in
\src{New/Reduction/Data/Existence}{the construction of sparse places}.

\paragraph{Two principal-level primes.}
All three level primes must be $1\bmod q_2^2$. We use the following
consequence of the classical Siegel--Walfisz theorem. For a fixed $C>0$
and all sufficiently large $W$, if $1\le Q\le W^C$ and
$W\le V\le W^C$, some prime $p\equiv1\bmod Q$ has
$V\le\log p\le V+W$. To see this, put $Z=e^{V+W}$. Then $Q\le(\log Z)^C$,
so the weighted prime count
$\sum_{p\le Z,\ p\equiv1\,(Q)}\log p$ has main term $Z/\varphi(Q)$
and error $O_C(Z e^{-c_C\sqrt{\log Z}})$. The error is
$o(Z/\varphi(Q))$ uniformly. Removing primes below $e^V$ costs
at most $Ve^V=o(Z/\varphi(Q))$, as $V,Q$ are polynomial in $W$.
We use the prime-weighted Siegel--Walfisz estimate of Montgomery and Vaughan \cite[Corollary~11.21]{MV}. Corollary~11.19 gives the von Mangoldt form. The readable edition works with
$Q=(q_0d)^{2(s+6)}$. Real characters at these levels are induced from the fixed modulus $(q_0d)^2$. The zero-free region and Newman's contour
argument then give \src{Lib/Analytic/Primes}{the restricted prime window}.
The classical estimate used here allows ineffective thresholds.

Let $h$ be the class number of $K$. For $i=0,\dots,h$ place logarithmic
windows of width $W$ at $10L+2iW$. They are disjoint and contained in
$[10L,11L]$ since $(2h+1)W\le L$ eventually. The preceding estimate
with $Q=q_2^2$ supplies $h+1$ distinct primes, all split in $K$
because they are $1\bmod d$. Select a prime ideal above each. Two have
the same ideal class. For their underlying primes $p_0,p_1$, the
product of the first ideal and the conjugate of the second is principal.
Its integral generator $\alpha$ has $\alpha\bar\alpha=M=p_0p_1$.
In particular $b^{20}\le M\le b^{22}$.

\paragraph{The level matched to the weights.}
Let $\kappa_*=\exp(f^{-1}\sum_{\tau\in\Phi}\log\kappa(\tau))$ and
$\delta_E$ be the root discriminant. After choosing $M$ put
\[
 V_*=m\log\kappa_*+(m-1)\log(M\delta_E/(2\pi))
                    -\tfrac{15}{2}L .
\]
We choose the third level prime using $V_*$, which now incorporates the mean weight of the selected places. Since
$E\subseteq\mathbb Q(\zeta_{q_1d})$, the root-discriminant tower
inequality and the cyclotomic discriminant formula give
$1\le\delta_E\le q_1d\le b$ for all large $x$. Also
$1\le\kappa_*\le b$. Hence
$19L\le\log(M\delta_E/(2\pi))\le23L$ eventually, and
\[
 mL\le V_*\le24mL
\]
since $m\ge2$. The modulus $q_2^2$ and $24mL$ are bounded by fixed
powers of $W$, and $V_*\ge W$. Apply the prime-window estimate once
more to choose a prime $N\equiv1\bmod q_2^2$ with
$V_*\le\log N\le V_*+W\le V_*+\sqrt L$.
Since $m\to\infty$, eventually $\log N\ge mL>11L$, so $N$ differs from $p_0$ and $p_1$. Finally $N\equiv1\bmod4$ is a sum of two integral squares and
$\sqrt{-1}\in E$. Thus $N=\beta\bar\beta$ with
$\beta\in\mathcal O_E$. This gives the required level and norm generators $\alpha$ and $\beta$. The bound $\log x_1\ge W-1\ge\sqrt W=x^{\eta/12}$ for large $W$ completes the proof of Theorem~\ref{thm:p1-data}.
The readable edition gives the \src{New/Reduction/Data/AdmissibleData}{definition of admissibility and
bounds on $V_*$} and
proves \src{New/Reduction/Data/Existence}{eventual existence with prescribed
integer length}.

\section{A one-sided prime number theorem}\label{sec:from-data}
Both secondary remainders in \eqref{eq:one} are evaluated at \(x_1\), with
\(q_0\) and \(d\) fixed as \(x\to\infty\). Since \(\log x_1\ge x^{\eta/12}\)
eventually, a sufficiently large inverse power of \(\log x_1\) absorbs both errors.

\begin{lemma}[Centering and variation]\label{lem:pnt}
Fix an odd prime $d$ and a positive integer $q$ coprime to $d$.
For every integer $K\ge1$ and all sufficiently large $y$,
\begin{equation}
 |R_{d,q}(y)|\le (\log y)^{-K}.\label{eq:PNT}
\end{equation}
For all $1\le u\le v$ one has
\begin{equation}
 |R_{d,q}(v)-R_{d,q}(u)|\le 13(v-u)/u.\label{eq:variation}
\end{equation}
\end{lemma}
\begin{proof}
Write $A_{d,q}(t)=\sum_{n\le t}\lambda_{d,q}(n)$ and $e(t)=A_{d,q}(t)-t/\varphi(q)$.
Integrating the step function gives
\begin{equation}
 R_{d,q}(y)=\int_1^y\frac{e(t)}{t^2}\,dt-c_{d,q}.\label{eq:integralR}
\end{equation}
Each \(n\le y\) contributes \(\lambda_{d,q}(n)\int_n^y t^{-2}\,dt\).
By \(0\le\lambda_{d,q}\le2\Lambda\) and Chebyshev's bound
\(\sum_{n\le t}\Lambda(n)\le6t\), we have \(|e(t)|\le13t\).
Integration over \([u,v]\) gives \eqref{eq:variation}.

For \eqref{eq:PNT}, apply the prime number theorem in progressions modulo
the fixed integer \(qd\)~\cite[Corollary~11.19]{MV}. On the reduced classes congruent to \(1\bmod q\),
the weights \(1-\chi_d\) sum to \(\varphi(d)\) by the Chinese remainder
theorem. Among nonunits modulo \(qd\), only powers of \(d\) can contribute,
with total weight \(O_d(\log(2t))\). Thus for each fixed \(B>1\),
\[
 A_{d,q}(t)=\frac{t}{\varphi(q)}+O_{q,d,B}\big(t/(\log t)^B\big)\qquad(t\ge2).
\]
Hence \(e(t)t^{-2}\) is integrable at infinity. For real \(s>1\), partial
summation and \eqref{eq:p1Dir} give
\[
 s\int_1^\infty e(t)t^{-s-1}\,dt
 =\sum_{n\ge1}\lambda_{d,q}(n)n^{-s}-\frac{s}{\varphi(q)(s-1)}
 =G_{d,q}(s).
\]
Letting \(s\downarrow1\) by dominated convergence gives
\(\int_1^\infty e(t)t^{-2}\,dt=G_{d,q}(1)=c_{d,q}\). Thus \(R_{d,q}(y)\)
is the negative of the tail integral. With \(B=K+2\) this is
\(O_{q,d,K}((\log y)^{-K-1})\), giving \eqref{eq:PNT} with coefficient one
for sufficiently large \(y\).
\end{proof}
The formal argument obtains \eqref{eq:PNT} from a quantitative Tauberian
theorem for the regularized logarithmic derivative. The fixed-modulus
prime number theorem used here follows from the classical zero-free estimate
near \(\Re s=1\).

\begin{theorem}[All real cutoffs]\label{thm:all-y}
Assume Hypothesis~\ref{hyp:scale}. For every $q_0$ divisible by $18480$, $q_0>0$, every prime $d\equiv q_0/2-1\pmod{q_0}$ and every $\varepsilon>0$ there is $Y=Y(q_0,d,\varepsilon)$ such that
\begin{equation}
 R_{d,q_0}(y)\le y^{-1/2+\varepsilon}\qquad(y\ge Y).\label{eq:real}
\end{equation}
\end{theorem}
\begin{proof}
Choose \(0<\eta<1/100\) with \(7\eta\le\varepsilon\) and an integer
\(K\ge24/\eta\). For every sufficiently large integer \(x\),
Hypothesis~\ref{hyp:scale} gives a datum of Theorem~\ref{thm:p1-data}
at \(x\) satisfying \eqref{eq:det}. On this datum \eqref{eq:p1-data-scale}
gives \(m\ge\sqrt{x}\), \(L\le2x^\eta\), and \(\log x_1\ge x^{\eta/12}\).
For \(q=q_0,1\), \eqref{eq:PNT} now gives \(|R_{d,q}(x_1)|\le x^{-2}\).
Substitution in \eqref{eq:one} yields
\begin{align*}
 R_{d,q_0}(x)&\le64\varphi(q_0)x^{-1/2+6\eta}
                 +(1+2x^\eta)x^{-2}\\
              &\le C x^t,\qquad C=64\varphi(q_0)+3,\quad t=-1/2+6\eta .
\end{align*}
The factor \(x^{-1/2}\) comes from \(m^{-1}\). For real \(y\ge2\) put
\(x=\lfloor y\rfloor\). Since \(x\ge y/2\) and \(-1\le t\le0\), we have
\(x^t\le2y^t\), and \eqref{eq:variation} bounds the increment by \(26/y\).
Hence for sufficiently large \(y\),
\[
 R_{d,q_0}(y)\le(2C+26)y^t\le y^{t+\eta}\le y^{-1/2+\varepsilon}.
\]
\end{proof}


\section{The transform and its real residues}\label{sec:p1polar}
Let \(d\) be an odd prime coprime to \(Q\). We use the entire functions
\(D_\eta\) and their regularized logarithmic derivative average \(G_{d,Q}\)
from \eqref{eq:p1G}. Under \(y=e^v\) the Mellin transform of the smoothed
remainder becomes a Laplace transform.
\begin{proposition}[Transform and residues]\label{prop:p1polar}
Put $E(v)=R_{d,Q}(e^v)$ and initially, for $\Re s>2$, set
\[
 F(s)=\int_0^\infty E(v)e^{(1-s)v}\,dv.
\]
Then
\begin{equation}\label{eq:p1F}
 F(s)=\frac{G_{d,Q}(s)/s-c_{d,Q}}{s-1}.
\end{equation}
Its meromorphic continuation has at most simple poles in $\Re s>0$ and is regular at $1$. Write
\[
 n_Q(z)=\sum_{\eta\bmod Q}\operatorname{ord}_z D_\eta,
 \qquad n_{Q,d}(z)=\sum_{\eta\bmod Q}\operatorname{ord}_z D_{\eta\chi_d}.
\]
With residue zero at regular points, for $\Re z>0$, $z\ne1$ one has
\begin{equation}\label{eq:p1res}
 \operatorname{Res}_{s=z}F(s)
     =\frac{n_Q(z)-n_{Q,d}(z)}{\varphi(Q) z(1-z)} .
\end{equation}
\end{proposition}
\begin{proof}
For the contribution of $n$ the integral starts at $\log n$. Direct integration gives
\[
 \int_{\log n}^\infty (n^{-1}-e^{-v})e^{(1-s)v}\,dv
       =\frac{n^{-s}}{s(s-1)}.
\]
The same integral at \(\Re s\) gives absolute convergence after summing
against \(\lambda_{d,Q}(n)\ge0\). The transforms of \(v/\varphi(Q)\) and
\(c_{d,Q}\) are \(1/[\varphi(Q)(s-1)^2]\) and \(c_{d,Q}/(s-1)\),
respectively. Substitution of \eqref{eq:p1Dir} cancels the double pole and
gives \eqref{eq:p1F}. Since \(c_{d,Q}=G_{d,Q}(1)\), the singularity at
\(1\) is removable. Elsewhere in \(\Re z>0\), each logarithmic derivative
contributes the order of its entire factor. The twisted orders enter
\(G_{d,Q}\) with positive sign, the untwisted with negative sign.
Dividing by \(z(z-1)\) gives \eqref{eq:p1res}.
\end{proof}
The remainder \(E\) is real and measurable. Since
\(0\le\lambda_{d,Q}(n)\le2\log n\), comparison with
\(\sum_{n\le y}2\log(n)/n\) bounds the smoothed sum by \(O(1+v^2)\)
at \(y=e^v\). Thus \(|E(v)|\ll_{d,Q}e^v\) for \(v\ge0\).
The corresponding formal constructions are
\src{Lib/Analytic/EntireLFunction}{the entire \(L\)-functions} and
\src{New/Reduction/Zeros/PrimeError}{the polar prime remainder}.

\subsection{Landau's argument on the real axis}
To deduce holomorphy from an upper bound on \(E\), we need the sign of the
residue at a possible real boundary pole. The following form of Landau's
singularity theorem makes the sign explicit \cite[Lemma~15.1]{MV}\cite[Chapter~II, Theorem~5b]{widder}.

\begin{lemma}[One-sided Landau criterion]\label{lem:p1landau}
Let $E$ be real measurable on $[0,\infty)$, with $|E(v)|\le Ce^v$. Suppose its transform $F$ as above continues meromorphically to $\Re s>0$ with at most simple poles there. Fix $\beta>0$. If $E(v)\le e^{(\beta-1)v}$ eventually and
\[
 \Re\operatorname{Res}_{s=u}F(s)\ge0
 \quad\text{for each real }u>\beta,
\]
then $F$ is holomorphic for $\Re s>\beta$.
\end{lemma}
\begin{proof}
The eventual upper bound and \(|E(v)|\le Ce^v\) on the initial bounded
interval give a constant \(K\ge1\) with \(E(v)e^v\le K e^{\beta v}\) for
every \(v\ge0\). Define the positive measure
\[
 d\mu(v)=(K e^{\beta v}-E(v)e^v)\,dv
 \qquad(v\ge0).
\]
Its Laplace transform $P(s)=\int e^{-sv}\,d\mu(v)$ converges for $\Re s>\max(2,\beta)$, where
\[
 P(s)=\frac K{s-\beta}-F(s).
\]
The right side continues \(P\) meromorphically to \(\Re s>0\).
Let \(\alpha\) be the abscissa of convergence of the positive integral,
allowing \(\alpha=-\infty\). Local absolute convergence makes \(P\)
holomorphic on \(\Re s>\alpha\).

Suppose $\alpha>\beta$. The continuation of $P$ has at most a simple pole at $\alpha$. Its residue has nonpositive real part by hypothesis, whereas
\[
 \lim_{u\downarrow\alpha}(u-\alpha)P(u)
\]
is a nonnegative real number because $P(u)$ is a positive integral there. The limit equals the residue, so the residue vanishes and the continuation is regular at $\alpha$.

Choose \(r>0\) so that the continuation is holomorphic on the disc of
radius \(3r\) about \(\alpha\), and set \(u=\alpha+r/2\). Differentiation
in the half-plane of convergence gives
\((-1)^kP^{(k)}(u)=\int v^k e^{-uv}\,d\mu(v)\).
The Taylor series at \(u\) converges at \(u-r<\alpha\). Its terms at this
point are nonnegative, so monotone convergence gives
\[
 \sum_{k\ge0}\frac{r^k}{k!}\int v^k e^{-uv}\,d\mu(v)
   =\int e^{-(u-r)v}\,d\mu(v).
\]
This contradicts the definition of $\alpha$. Hence $\alpha\le\beta$, and uniqueness of meromorphic continuation identifies $F=K/(s-\beta)-P$ throughout $\Re s>\beta$. Both terms are holomorphic there.
\end{proof}
The twist in the next section must therefore control residues on the real
axis as well as the residue at the proposed zero.

\section{A quadratic twist free of zeros in a prescribed box}
\label{sec:p1-good-twist}

Let $Q$ be a positive multiple of $8$ and $a_0\bmod Q$ a unit class with
$a_0\equiv7\pmod8$. For each prime $d\equiv a_0\pmod Q$ put
$\chi_d(n)=(n/d)$. We shall choose $d$ after fixing the height $H$ and
left edge $\beta$ of the box containing the proposed zero.

\begin{proposition}\label{prop:p1-good-twist}
For $H\ge0$ and $\beta>1/2$ there is a prime $d\equiv a_0\pmod Q$ such that
\[
 L(s,\xi\chi_d)\ne0
 \qquad (\xi\bmod Q,\ \Re s\ge\beta,\ |\Im s|\le H).
\]
\end{proposition}

It suffices to treat $1/2<\beta<1$, since the box with $\beta=3/4$
contains every box of the same height with $\beta\ge1$. Put
$\delta=(\beta-1/2)/9$. A zero forces a large theta integral. Its mean
over quadratic symbols bounds the number of exceptional primes.

\subsection{Detection by a truncated Mellin transform}
For $j=0,1$ and a sequence $c(n)$ of modulus at most one, put
\[
 w_j(y)=y^j e^{-\pi y^2},\qquad
 \theta_{j,c}(u)=u^{1/2}\sum_{n\ge1}c(n)w_j(nu),\qquad
 G_j(s)=\tfrac12\pi^{-(s+j)/2}\Gamma((s+j)/2).
\]
Thus $G_j$ is the Mellin transform of $w_j$. Let $\chi$ be primitive and
nonprincipal of conductor $q$ and parity $j$. Suppress $j$ in $\theta$,
write $\mu$ for the M\"obius function, and for $T\ge\sqrt q$ put
\[
 J(\chi,T)=\int_{1/T}^{\infty}
       (|\theta_\chi(u)|+|\theta_{\mu\chi}(u)|)\frac{du}{u},\qquad
 M_T(s)=\int_{1/T}^{\infty}\theta_{\mu\chi}(u)u^{s-1/2}\frac{du}{u}.
\]
\begin{lemma}\label{lem:p1-twist-detect}
For $T\ge\sqrt q$ above a threshold depending only on $\delta,H$, a zero $s$ of
$L(s,\chi)$ in the box of Proposition~\ref{prop:p1-good-twist} implies
$J(\chi,T)>T^{2\delta}$.
\end{lemma}
\begin{proof}
For $\Re s>1$, termwise integration gives the Mellin transform
$G_j(s)L(s,\chi)$ of $\theta_\chi$ at $s-1/2$. Poisson summation for
the primitive character gives
$|\theta_\chi(u)|=|\theta_{\bar\chi}(1/(qu))|$.
The weight $w_j$ is real, so $|\theta_{\bar\chi}|=|\theta_\chi|$.
Gaussian decay at infinity therefore gives rapid decay at zero and
continues the Mellin integral to every $s$. Splitting at $q^{-1/2}$
and substituting $u\mapsto1/(qu)$ in the lower integral gives,
on $\Re s=1/2$,
\[
 |G_j(s)L(s,\chi)|\le 2\int_{q^{-1/2}}^\infty|\theta_\chi(u)|\frac{du}{u}
 \le2J(\chi,T),\qquad |M_T(s)|\le J(\chi,T).
\]
On $\Re s=2$, absolute convergence identifies the full Mellin transform
for $\mu\chi$ with $G_j(s)/L(s,\chi)$. For $0<u\le1$ the Gaussian
sum gives $|\theta_{\mu\chi}(u)|\ll u^{-1/2}$, so truncation at $1/T$
costs $O(T^{-1})$ uniformly in $\chi$ and $\Im s$. The Dirichlet
series bounds $L(2+it,\chi)$ uniformly, and the Mellin integral for
$w_j$ bounds $G_j$ throughout $1/2\le\Re s\le2$. Thus
\[
 F_T(s)=G_j(s)L(s,\chi)M_T(s)-G_j(s)^2
\]
is $O(T^{-1})$ on $\Re s=2$. If $J(\chi,T)\le T^{2\delta}$, the
bounds on $\Re s=1/2$ give $F_T(s)\ll T^{4\delta}$ on that line. For fixed
$\chi,T$ the Mellin integrals give a bounded holomorphic function on
the strip, continuous on its boundary. The three-lines theorem~\cite[Lemma~4.6]{brown} now gives,
for $\beta\le\Re s\le1$,
\[
 |F_T(s)|\ll T^{4\delta(1-\omega)-\omega}\le T^{-2\delta},\qquad
 \omega=\frac{\Re s-1/2}{3/2}\ge6\delta .
\]
At a zero of $L(s,\chi)$ we have $|G_j(s)|^2\ll T^{-2\delta}$.
For $j=0,1$, the gamma factor has a positive lower bound in modulus
on $1/2\le\Re s\le1$, $|\Im s|\le H$, a contradiction for large $T$.
The Euler product excludes zeros for $\Re s>1$.
\end{proof}

\subsection{Mean over quadratic symbols}
For $n\ge1$ let $n^\circ$ be its odd part and put
$b_v(n)=\bigl(\frac{-v}{n^\circ}\bigr)$, the Jacobi symbol.
For either parity, any $|\alpha(n)|\le1$, and $u\ge1/T$,
\begin{equation}\label{eq:p1-twist-second}
 \sum_{1\le v\le T^2}|\theta_{j,\alpha b_v}(u)|^2
 \ll_{\delta} T^2u^{-2\delta}.
\end{equation}
Gaussian decay permits termwise expansion at fixed $u$. For a pair
$n,m$ the character sum is $\sum_{v\le T^2}b_v(nm)$. When $(nm)^\circ$
is square, its modulus is at most $T^2$. Using $|w_j(y)|\ll_A y^{-A}$
with $A=1/2+\delta$ bounds the pair's contribution by
$O_\delta(T^2u^{-2\delta}(nm)^{-1/2-\delta})$. Write
$n=2^h t k^2$ and $m=2^i t l^2$, with $t$ odd squarefree and $k,l$ odd.
Summation in $h,i$ gives geometric series. The exponents of $t,k,l$
are all $-1-2\delta$, so the remaining sums converge.

When $(nm)^\circ$ is nonsquare, $b_v(nm)$ is a fixed sign times a
nonprincipal character of $v$ modulo $(nm)^\circ$. P\'olya--Vinogradov
\cite[Theorem~9.18]{MV} gives an interval sum
$\ll ((nm)^\circ)^{1/2}(1+\log((nm)^\circ))\ll_\delta(nm)^{1/2+\delta/2}$.
Taking $A=3/2+\delta$ gives
$O_\delta(u^{-2-2\delta}(nm)^{-1-\delta/2})$ per pair. Summing in
$n,m$ and using $u^{-2}\le T^2$ proves \eqref{eq:p1-twist-second}.

Write $J_j(c,T)$ for the integral defining $J$ with coefficients $c$
and parity $j$. Cauchy--Schwarz in $v$ applied to
\eqref{eq:p1-twist-second} bounds the sum of absolute values by
$O_\delta(T^2u^{-\delta})$ for both $\alpha$ and $\mu\alpha$.
Integration gives
\begin{equation}\label{eq:p1-twist-mean}
 \sum_{v\le T^2,\ v\ge1} J_j(\alpha b_v,T)
 \ll_\delta T^2\int_{1/T}^\infty u^{-\delta-1}\,du
 \ll_\delta T^{2+\delta}.
\end{equation}

\begin{proof}[Proof of Proposition~\ref{prop:p1-good-twist}]
Let $\psi$ be the primitive character of conductor $M\mid Q$ inducing
$\xi\bmod Q$. Every prime $d\equiv a_0\pmod Q$ is coprime to $Q$,
so $\psi\chi_d$ is primitive of conductor $Md$ and nonprincipal. The
finite Euler factors relating $L(s,\xi\chi_d)$ to $L(s,\psi\chi_d)$
are nonzero for $\Re s>0$. Since $d\equiv7\pmod8$, we have
$\chi_d(2)=1$. Quadratic reciprocity gives $\chi_d(n)=b_d(n)$ for
all $n\ge1$, including multiples of $d$. The parity $j$ of $\psi\chi_d$
is opposite to that of $\psi$ for every $d$ in the progression.

Take $X$ sufficiently large and set $T=\sqrt{QX}$. For $d\le X$ the
conductor $Md$ is at most $T^2$. A zero of $L(s,\xi\chi_d)$ in the box
therefore gives $J_j(\psi b_d,T)>T^{2\delta}$ by
Lemma~\ref{lem:p1-twist-detect}. Equation~\eqref{eq:p1-twist-mean}
bounds the number of such primes by $O_{\delta,H,Q}(T^{2-\delta})$.
Let $B(X)$ count primes $d\le X$ in the progression for which some
$\xi\bmod Q$ has a twist with a zero in the box. Summing over the
$\varphi(Q)$ characters gives
\[
 B(X)\ll_{Q,H,\delta} X^{1-\delta/2}.
\]
Dyadic summation shows that the reciprocals of the exceptional primes
have finite sum. By Dirichlet's theorem in its analytic form~\cite[Chapter~4]{MV}, the
reciprocals of the primes in the reduced class have divergent sum.
Infinitely many primes in the class satisfy the proposition.
\end{proof}

The formal argument uses the weight $(\log d)/d$ in the last step.
The exceptional sum still converges by the same dyadic bound, and the
sum over the progression diverges. The Gaussian normalization in the
Poisson transformation, with Gauss sum of absolute value $\sqrt q$,
is given in \src{Lib/Analytic/DirichletTheta}{Dirichlet theta transforms}.
The weighted selection argument is in
\src{Lib/Analytic/QuadraticTwist}{quadratic twist selection}.

\section{Nonvanishing and reflection}\label{sec:descent}
\subsection{Choosing the twist at the height of the proposed zero}
\begin{theorem}[Family criterion]\label{thm:p1criterion}
Let $8\mid Q$, $Q>0$, and let $d_0$ be prime to $Q$ with $d_0\equiv7\bmod 8$. Fix $\beta>1/2$. Suppose that for each prime $d\equiv d_0\bmod Q$ there exists $Y_d$ such that
\[
 R_{d,Q}(y)\le y^{\beta-1}\qquad(y\ge Y_d).
\]
Then for every character $\psi$ modulo $Q$ and every $\rho\ne1$ with $\Re\rho>\beta$, one has $L(\rho,\psi)\ne0$.
\end{theorem}
\begin{proof}
Fix $\rho$ with $\Re\rho>\beta$, $\rho\ne1$. Apply Proposition~\ref{prop:p1-good-twist} with $H=|\Im\rho|$. It gives a prime $d\equiv d_0\bmod Q$ for which $n_{Q,d}$ vanishes throughout $\Re s\ge\beta$, $|\Im s|\le H$. In particular $(d,Q)=1$ and $d\equiv7\bmod 8$. For this $d$ the assumed bound gives $E(v)=R_{d,Q}(e^v)\le e^{(\beta-1)v}$ whenever $e^v\ge Y_d$.

For real $\beta<u<1$, formula~\eqref{eq:p1res} now gives residue $n_Q(u)/[\varphi(Q)u(1-u)]\ge0$. The transform is regular at $1$. For $u>\max(1,\beta)$ the untwisted factors are also nonzero, so the residue is zero. Lemma~\ref{lem:p1landau} therefore gives holomorphy in $\Re s>\beta$. At $\rho$ the residue vanishes, and $n_{Q,d}(\rho)=0$ by the choice of $H$. Thus \eqref{eq:p1res} gives $n_Q(\rho)=0$. This is a sum of nonnegative orders, so $D_\psi(\rho)\ne0$ for every $\psi\bmod Q$. As $\rho\ne1$, the same holds for $L(\rho,\psi)$.
\end{proof}
In the \src{New/Reduction/Zeros/Criterion}{one-sided nonvanishing theorem}, the height $H$ determines the box for both the residue at $\rho$ and the signs on the real axis. The threshold $Y_d$ may depend on the prime chosen for this box.

\subsection{Return to the original modulus}
\begin{proof}[Proof of Theorem~\ref{thm:main}]
Fix a character $\chi\bmod q$, $q\ge1$. On $\Re s\ge1$ classical nonvanishing gives the assertion, with the principal pole at $1$ excepted. Take $\rho$ with $1/2<\Re\rho<1$ and choose $1/2<\beta<\Re\rho$. Put
\[
 Q=18480q,\qquad d_0=Q/2-1=9240q-1 .
\]
Here $8\mid Q$, $d_0\equiv7\bmod 8$, and $(d_0,Q)=1$ since $d_0$ is odd and $Q=2(d_0+1)$. Under Hypothesis~\ref{hyp:scale}, Theorem~\ref{thm:all-y} with $\varepsilon=\beta-1/2$ gives $R_{d,Q}(y)\le y^{\beta-1}$ for all sufficiently large $y$, for each prime $d\equiv d_0\bmod Q$. Apply Theorem~\ref{thm:p1criterion} to the character $\chi_Q$ induced at level $Q$. It gives $L(\rho,\chi_Q)\ne0$. The identity
\begin{equation}
 L(s,\chi_Q)=L(s,\chi)\!
 \prod_{\substack{p\mid Q\\p\nmid q}}(1-\chi(p)p^{-s})\label{eq:level}
\end{equation}
holds by the Euler product for $\Re s>1$ and elsewhere by meromorphic continuation. At $s=\rho$ the left side is nonzero, so $L(\rho,\chi)\ne0$.

Let $\psi$ be the primitive character inducing $\chi$. The additional Euler factors in \eqref{eq:level}, applied from the conductor of $\psi$ to $q$, are nonzero for $\Re s>0$ because $|\psi(p)p^{-s}|<1$. A zero of $L(s,\chi)$ in the open critical strip is therefore a zero of $L(s,\psi)$. The gamma factors are finite and nonzero in the strip, and the primitive functional equation~\cite[\S10.1]{MV} gives $L(1-s,\bar\psi)=0$. Half-plane nonvanishing at $s$ and $1-s$ gives $\Re s\le1/2$ and $1-\Re s\le1/2$, respectively.

For $q=1$ this gives the half-plane assertion for $\zeta$. Suppose $\zeta(s)=0$ outside the negative even integers. Then $s\ne1$ and $s\ne0$, since $\zeta(0)=-1/2$. The factor $\pi^{-s/2}\Gamma(s/2)$ is finite and nonzero at $s$, so the completed zeta function vanishes at both $s$ and $1-s$. If $\Re s<1/2$, then $\Re(1-s)>1/2$ and $1-s\ne1$. The gamma factor is finite and nonzero there, so $\zeta(1-s)=0$, contrary to half-plane nonvanishing. The same assertion rules out $\Re s>1/2$.
\end{proof}

\subsection{The converse on the constructed scales}
\begin{corollary}[Strength of the determinant criterion]\label{cor:scale-equivalence}
Hypothesis~\ref{hyp:scale} is equivalent to GRH for Dirichlet \(L\)-functions.
Under GRH, for fixed permissible \((q_0,d)\) and fixed \(0<\eta<1/100\),
every admissible datum at the scales of \eqref{eq:p1-data-scale} satisfies
\(m|\Tot_D|\le fa(\log b)^5\) at every sufficiently large integer \(x\).
\end{corollary}
\begin{proof}
Theorem~\ref{thm:main} gives the forward implication.
Assume GRH and take \(Q=1\) or \(Q=q_0\). Write
\(\psi(t,\xi)=\sum_{n\le t}\Lambda(n)\xi(n)\). Orthogonality gives
\[
 A_{d,Q}(t)=\frac1{\varphi(Q)}
  \sum_{\xi\bmod Q}\bigl(\psi(t,\xi)-\psi(t,\xi\chi_d)\bigr).
\]
There is one principal character in the first sum. Every twist in the
second sum is nonprincipal by the Chinese remainder theorem, including
when \(Q=1\). The prime number theorem under GRH at the fixed levels
\(Q,Qd\) gives \(A_{d,Q}(t)-t/\varphi(Q)\ll_{d,Q}\sqrt{t}\log^2(2t)\)
\cite[Theorem~13.7]{MV}. Imprimitive Euler factors change the corresponding
primitive sums by \(O_{d,Q}(\log(2t))\).
By \eqref{eq:integralR} and Lemma~\ref{lem:pnt}, \(R_{d,Q}(y)\) equals
the negative tail integral of this centered error against \(dt/t^2\).
Thus \(\lvert R_{d,Q}(y)\rvert\ll_{d,Q} y^{-1/2}\log^2(2y)\).

The fibre of $-1\bmod q_0$ has $A$ members, so its selected subset has $B\le A$. Substitution in \eqref{eq:count}, with $f=A\varphi(q_0)$ and $m<2\sqrt{x}$, gives
\[
 \frac{m|\Tot_D|}{fa}\ll_{d,q_0}\log^2(2x)
       +\sqrt{x/x_1}\log^2(2x_1)=o(L^5).
\]
Indeed \(x_1=\lfloor\exp(x^{\eta/6})\rfloor\) makes the second term tend
to zero, and \(L>x^\eta\). The estimate is uniform over the data at the
fixed parameters. Theorem~\ref{thm:p1-data} supplies a datum at each
sufficiently large integer \(x\), proving Hypothesis~\ref{hyp:scale}.
\end{proof}
